\subsection{The diffuse kernel outside the normalizer algebra}
We retain the faithful normal state $\varphi$, with $A\subset M_\varphi$,
and the algebra $N$ generated by the groupoid normalizers of $A$.
Use a standard probability model $A=L^\infty(V,\mu)$ for
$\varphi|_A$.  We use the standard-form realization $L^2(M)$, equivalently
Haagerup $L^2(M)$, and write $\varphi^{1/2}$ for the positive vector
implementing $\varphi$.  In particular,
$a\varphi^{1/2}=\varphi^{1/2}a$ for $a\in A$.

For $X\in M$, the commuting left and right actions of $A$ give a finite
joint spectral measure $\nu_X$ for $X\varphi^{1/2}$ on $V\times V$.
Our first coordinate is the right coordinate: for Borel sets $B,C\subset V$,
\begin{equation}\label{exp:joint-measure}
 \nu_X(B\times C)
 =\big\|1_CX\varphi^{1/2}1_B\big\|_2^2
 =\int_B E_A(X^*1_CX)(v)\,d\mu(v).
\end{equation}
Both marginal measures are absolutely continuous with respect to $\mu$.
The right marginal has density $E_A(X^*X)\leq\|X\|^2$; the left marginal
is absolutely continuous because $a\mapsto\varphi(X^*aX)$ is normal on
$A$.  Disintegration on standard Borel spaces therefore gives a measurable
kernel $k_X(v,dw)$ such that
\begin{equation}\label{exp:kernel}
 \nu_X(dv,dw)=d\mu(v)k_X(v,dw),\qquad
 E_A(X^*aX)(v)=\int_V a(w)\,k_X(v,dw),\qquad
 k_X(v,V)\leq\|X\|^2
\end{equation}
almost everywhere, for $a\in A$.  As usual, individual function identities
are understood modulo null sets.  The joint spectral resolution can be
constructed from a bounded real coordinate generating the standard model
and the commuting left and right spectral resolutions of that coordinate.

\begin{lemma}\label{lem:atomless}
 If $E_N(X)=0$, then $k_X(v,\cdot)$ has no atoms for almost every $v\in V$.
\end{lemma}

\begin{proof}
Choose Borel representatives for the kernel.  The set
\[
 \mathcal T=\{(v,w):k_X(v,\{w\})>0\}
\]
is Borel and has countable sections in the second coordinate.  For example,
joint measurability of the singleton mass follows by using shrinking
neighborhoods in a compatible separable metric on the standard Borel
space.  The countable-section uniformization theorem expresses
$\mathcal T$ as a countable union of graphs of partial Borel maps.
If the conclusion fails, one such map $F:S\to V$ has a domain of positive
$\mu$-measure, with
\[
 g(v):=k_X(v,\{F(v)\})>0\quad(v\in S).
\]
There is no assumption that $F$ is injective.  For a $\mu$-null Borel set
$B\subset V$, absolute continuity of the left marginal gives
\[
 \int_{S\cap F^{-1}(B)}g(v)\,d\mu(v)
 \leq\nu_X(V\times B)=0.
\]
Thus $F$ pulls back null sets to null sets.  With $e=1_S\in A$, composition
with $F$ consequently defines a normal $*$-homomorphism
\[
 \theta:A\longrightarrow Ae,\qquad
 \theta(a)(v)=a(F(v))\quad(v\in S),\qquad \theta(1)=e.
\]
The graph will intertwine $A$ with the possibly proper subalgebra
$\theta(A)$ of $Ae$; its polar partial isometry need not yet normalize
$A$. We will use an expected type I intermediate algebra to split the
right support into corners where these two abelian algebras coincide.

Apply the joint spectral projection of the graph of $F$ to
$X\varphi^{1/2}$ and call the resulting vector $\zeta$.  It is nonzero and
satisfies
\begin{equation}\label{exp:graph-intertwining}
 \zeta e=\zeta,\qquad a\zeta=\zeta\theta(a)\quad(a\in A).
\end{equation}
Its right spectral density on $S$ is $g$.  In particular, $e$ is its
smallest right supporting projection among the projections of $A$:
if $f\in A$ is a projection and $\zeta f=\zeta$, then $e\leq f$.

Take the $L^2$ polar decomposition
$\zeta=u h$, where $h=|\zeta|\in L^2(M)_+$ and $u\in M$ is a partial
isometry.  Put $p=u^*u\leq e$ and $q=uu^*$.  We use here the usual
polar decomposition and support properties of standard-form
noncommutative $L^2$; see, for example, \cite[Lemma~11.12 and Theorem~11.30]{Hiai20}.
For a unitary $a\in A$, the operator $b=\theta(a)+1-e$ is unitary and
\eqref{exp:graph-intertwining} reads $a\zeta=\zeta b$.
Taking absolute values gives $h=b^*hb$.  Uniqueness of the polar
decomposition then gives $au=ub=u\theta(a)$.  Consequently,
\begin{equation}\label{exp:polar-intertwining}
 h\theta(a)=\theta(a)h,\qquad
 p\theta(a)=\theta(a)p,\qquad
 au=u\theta(a)\quad(a\in A),\qquad q\in A.
\end{equation}
The identities for all $a$ follow by linearity from unitaries.  For the
last assertion, the identities imply that $q$ commutes with $A$, and $A$
is maximal abelian.

Consider the von Neumann algebra
\[
 D=\theta(A)'\cap eMe.
\]
It contains $p$ and $Ae$, and $Ae$ is a MASA in $D$.
Moreover $\theta(A)\subset Z(D)$.  If $d\in pDp$, then $udu^*$ commutes
with $Aq$, by \eqref{exp:polar-intertwining}, and hence belongs to $Aq$.
On the other hand,
\[
 u\theta(a)p u^*=aq\quad(a\in A).
\]
Conjugation by $u$ therefore identifies $pDp$ with $Aq$, and
\begin{equation}\label{exp:abelian-corner}
 pDp=\theta(A)p.
\end{equation}
Let $c$ be the central support of $p$ in $D$.  Since
$Z(D)\subset Ae$, the identity $pc=p$ gives $\zeta c=\zeta$; the minimal
right $A$-support property gives $c=e$.  Thus $D$ has an abelian projection
with full central support and is type I.  Finally, if $d\in Z(D)$, then
\eqref{exp:abelian-corner} gives $dp=\theta(a)p$ for some $a\in A$.
The central element $d-\theta(a)$ annihilates a projection of full central
support and is zero.  We have proved
\begin{equation}\label{exp:center}
 Z(D)=\theta(A).
\end{equation}

We next need a partition of $e$ by projections $e_j\in Ae$ such that
\begin{equation}\label{exp:branch-partition}
 (Ae)e_j=\theta(A)e_j.
\end{equation}
The restriction of $E_A$ to $D$ is a normal conditional expectation onto
$Ae$.  The type I MASA theorem of Akemann and Sherman
\cite[Theorem~4.1]{AS12} supplies a partition of $e$ by projections in
$Ae$ that are abelian in $D$.  There are only countably many nonzero
members, because $D$ has the faithful normal state
$\varphi|_D/\varphi(e)$.  For an abelian projection in a type I algebra,
\[
 e_jDe_j=Z(D)e_j.
\]
This follows, for instance, by writing the type I algebra as a direct
integral: an abelian projection has fiber rank at most one.  Together
with \eqref{exp:center} and $\theta(A)\subset Ae\subset D$, this proves
\eqref{exp:branch-partition}.

Set $\eta_j=\zeta e_j$ and take its polar decomposition
$\eta_j=u_j h_j$.  Its right support is at most $e_j$.  Repeating
\eqref{exp:polar-intertwining} shows that $h_j$ commutes with $\theta(A)$
and that
\[
 au_j=u_j\theta(a)\quad(a\in A).
\]
By \eqref{exp:branch-partition} and the support condition, $h_j$ commutes
with all of $A$.  Therefore the right support $p_j=u_j^*u_j$ belongs to
$A$; the left support $q_j=u_ju_j^*$ belongs to $A$ by the same polar
argument as before.  In addition,
\[
 \theta(A)p_j=Ap_j,\qquad
 u_j(Ap_j)u_j^*=Aq_j.
\]
Thus each nonzero $u_j$ is a groupoid normalizer and belongs to $N$.

Let $\psi_j$ be the normal positive functional whose positive
implementing vector is $h_j$.  Since $h_j$ commutes with $A$, this
functional centralizes $A$.  In particular
$\psi_j\circ C_{\mathcal Q}=\psi_j$ for every finite partition
$\mathcal Q$ in $A$.  Lemma~\ref{lem:compression-net}, normality, and
boundedness of the compression net imply
\[
 \psi_j=\psi_j\circ E_A.
\]
Write $f_j=d(\psi_j|_A)/d\mu\in L^1(V,\mu)_+$ and
$f_{j,m}=\min(f_j,m)$.  The positive vectors
$f_{j,m}^{1/2}\varphi^{1/2}$ belong to $A\varphi^{1/2}$ and converge in
$L^2$; indeed their squared pairwise distances are
\[
 \int_V\big|f_{j,m}^{1/2}-f_{j,n}^{1/2}\big|^2\,d\mu.
\]
Their normal positive functionals are
$x\mapsto\int f_{j,m}E_A(x)\,d\mu$ and converge in norm to $\psi_j$.
The limit is therefore the unique positive implementing vector $h_j$.
It follows that
\[
 h_j\in\overline{A\varphi^{1/2}},\qquad
 \eta_j\in\mathcal H_N:=\overline{N\varphi^{1/2}}.
\]
The right projections $e_j$ are orthogonal and sum to $e$, so
$\zeta=\sum_j\zeta e_j\in\mathcal H_N$, with convergence in $L^2$.

The closed space $\mathcal H_N$ reduces both the left and right actions
of $A$: for the right action, use
$n\varphi^{1/2}a=na\varphi^{1/2}$.  It therefore reduces their joint
spectral projections, including the graph projection defining $\zeta$.
But $E_N(X)=0$ gives, for $n\in N$,
\[
 \varphi(n^*X)=\varphi(n^*E_N(X))=0,
\]
so $X\varphi^{1/2}\perp\mathcal H_N$.  Thus
$\zeta\perp\mathcal H_N$, contradicting $\zeta\ne0$ and the
membership just proved.
\end{proof}

\paragraph{Removing atomic summands.}
For the remaining treatment of $Z=x-E_N(x)$ we may assume that $A$ is
diffuse.  Here are the details of this reduction.  If $f$ is an atom of
$A$, then $Af=\mathbb C f$ is a MASA in $fMf$, so $fMf=\mathbb C f$.
Thus $f$ is minimal in $M$.  Let $z_f$ be its central support in $M$.
The algebra $z_fM$ is type I, since $f$ is an abelian projection with
full central support there.  Compression by $f$ is faithful on its
center and has scalar range, so $z_fM$ is a type I factor.

The algebra $Az_f$ is a MASA in this factor and has the normal
expectation $E_A|_{z_fM}$.  Again by \cite[Theorem~4.1]{AS12}, it has a
partition of $z_f$ into rank-one projections.  In particular it is
purely atomic.  Partial isometries joining any two of these rank-one
projections normalize their $A$-corners.  A compatible system of these
matrix units generates $z_fM$, and hence
\[
 z_fM\subset N.
\]
Distinct central supports of atoms are either equal or orthogonal,
because each supports a factor summand.  Their sum $z_{\mathrm{at}}$ is
a central projection, belongs to $A$, and satisfies
$z_{\mathrm{at}}M\subset N$.  Consequently
\[
 z_{\mathrm{at}}Z=0.
\]
On the complementary central summand, $A(1-z_{\mathrm{at}})$ has no
atoms.  The algebra generated there by its groupoid normalizers is
$N(1-z_{\mathrm{at}})$: corner normalizers extend by zero, and global
normalizers restrict to this central summand.  We may use the normalized
restriction of $\varphi$ on the nonzero complementary summand.  This
changes state seminorms only by a fixed scalar factor.
A partition obtained there extends to $M$ with the same number of cells
by adding $z_{\mathrm{at}}$ to one cell; all estimates for $Z$ are
preserved because $Z$ vanishes on $z_{\mathrm{at}}$.
